\section*{A Appendix}
We will prove Theorem9, which we state again here for convenience.
\paragraph{Theorem32 (Theorem9).} Let $m=10n$. Then there exists constants $k>0$, $0<\epsilon<1$ such that for sufficiently large $n$, there exists $k$-CSP formulas $\{F_{n,k}\}$ with $m$ $k$-local constraints such that for $m'=(1-\epsilon)m$, every subset of $m'$ constraints is unsatisfiable and requires linear degree PC refutations.
First we'll show that a random regular bipartite graph has good boundary expansion. This has been used implicitly in other works ([9],[5]), but for completeness we state and prove it here. Let $G=(L,R,E)$ be a bipartite graph, and let $A\subseteq R$. The boundary for $A$, $\partial(A)$, is the set of vertices $x$ in $L$ so that $|N(x)\cap A|=1$, i.e., vertices with a unique neighbor in $A$. A bipartite graph is $(d,k)$ regular if every vertex in $L$ has degree $d$ and every vertex in $R$ has degree $k$. In this case, for $n=|L|$, $m=|R|$, we have $dn=km$.
\paragraph{Theorem33.} Let $d,k,n,m$ be positive integers with $dn=km$, $k\geq12$. Then with high probability for a random $(d,k)$ regular bipartite graph with $|L|=n$, $|R|=m$, for all $A\subset R$, $|A|<n/(e^6k^2)$, we have $\partial(A)\geq k|A|/2$.
\begin{proof}
Let $N(A)$ be all the neighbors of $A$. Since the total degrees of vertices in $A$ is $k|A|$, and each element of $N(A)-\partial(A)$ is contingent on two such edges, $k|A|\geq2(|N(A)|-|\partial(A)|)+|\partial(A)$, or $\partial(A)\geq2|N(A)|-k|A|$. We will show that with high probability for all such $A$, $|N(A)|>3k|A|/4$, and hence $\partial(A)\geq k|A|/2$.
If not, there are sets $A\subset R$ and $B\subset L$ so that $N(A)\subseteq B$ and $|B|=3k|A|/4$. We will bound the probability that this is true for fixed sets $A,B$ and then take a union bound. We can view picking a random $(d,k)$ bipartite graph as picking a random matching between $d$ half-edges adjacent to each $x\in L$ and $k$ such half-edges adjacent to each $y\in R$; if a half edge for $x$ is matched to a half-edge for $y$, it forms an edge between $x$ and $y$.
We can form this matching by going through the half edges for nodes in $R$ and for each randomly selecting an unmatched half-edge for some node in $L$. We start with the edges for $A$ in an arbitrary order. If we condition on all previous neighbors for $A$ being in $B$, the number of half-edges left still available for $B$ is less than $d|B|$, whereas the number for $\bar B$ stays at exactly $d(n-|B|)$. Thus, the conditional probability that the next edge formed is also in $B$ is at most $|B|/n$, and we do this for each of $k|A|$ edges, meaning the probability that all neighbors are in $B$ is at most $(|B|/n)^{k|A|}$.
Now, for a fixed $|A|$ and setting $|B|=3k|A|/4$, we take the union bound over all subsets $A$ and $B$. This gives a total probability of failure for some set $A$ of size $a$ as:
\begin{align*}
&\binom ma\binom n{3ka/4}(3ka/4n)^{ka}\\
&\leq(em/a)^a(4en/3ka)^{3ka/4}(3ka/4n)^{ka}\\
&\leq(em/a)^a(e^3ka/n)^{ka/4}\\
&=(ekn/da)^a(e^3ka/n)^{ka/4}\\
&=\left(\frac{e^{3k/4+1}a^{k/4-1}k^{k/4+1}}{dn^{k/4-1}}\right)^a.
\end{align*}
Since we are assuming $a<n/(e^6k^2)$, the base in the above expression is at most
\[\frac{e^{3k/4+1}(n/e^6k^2)^{k/4-1}k^{k/4+1}}{dn^{k/4-1}}=e^{7-3k/4}k^{3-k/4}/d\]
which for $k\geq12$ is bounded below $e^{-2}$, meaning the probability of such a bad set existing is exponentially small in $a$, and the probability of such a bad set existing for any $a$ is less than $1/2$.
\end{proof}
\paragraph{Definition34.} For a Boolean vector $X=\{x_1,\ldots,x_n\}$, we define $\mathcal L_{n,m,k_1,k}(X)$ to be the distribution over $k$-CSP formulas over $n$ variables $X=\{x_1,\ldots,x_n\}$ obtained by selecting $m$ parity equations, where each parity is represented by a node on the right of a randomly chosen bipartite graph $G(L,R,E)$, with $|L|=n$, $|R|=m$, and with left degree bounded by $k_1$ and right degree bounded by $k$.
\paragraph{Lemma35.} Let $F_{n,k}$ be a tautology given by the system of parity equations $AX=b$ over variables $X=\{x_1,\ldots,x_n\}$ drawn at random from $\mathcal L_{n,m,k_1,k}$ where $m=10n$, for large enough constants $k_1,k>0$, and $b$ is chosen randomly. Then the following hold with high probability for a small enough $\epsilon>0$:
\begin{enumerate}[label=\alph*)]
\item Any subset of a $(1-\epsilon)$-fraction of the equations in $F_{n,k}$ is unsatisfiable.
\item Any subset of a $(1-\epsilon)$-fraction of the equations in $F_{n,k}$ requires PC degree $c_2n$ to refute, for some $c_2>0$.
\end{enumerate}
\begin{proof}
\begin{enumerate}[label=\alph*)]
\item The probability that a set of $(1-\epsilon)10n$ random parities (i.e. for a random choice of $b$) is satisfiable is at most $2^{-9n}$ for a small enough $\epsilon$. The probability that any such subset of $F_{n,k}$ is satisfiable is therefore at most $2^{(-n(9-10H(\epsilon)))}$, which is exponentially small for a small enough $\epsilon$ (where $H(\epsilon)$ is the binary entropy function).
\item This follows directly from [1], Theorem3.8 and Theorem4.4, since by Theorem33 the bipartite graph underlying the system of parity equations $A$ has good boundary expansion with high probability.
\end{enumerate}
\end{proof}
